\section{Introduction}
\label{sec:introduction}

Online sexism is widespread, and automated detection is one of the tools
platforms and researchers use to make identifying it tractable at scale
\cite{kirk2023semeval}. The Explainable Detection of Online Sexism
(EDOS) benchmark, which this paper evaluates against throughout, was
built specifically to move that detection effort past a coarse
sexist/not-sexist label toward a fine-grained taxonomy of how sexism
shows up in text. Most computational work on EDOS, and on hate speech
detection more broadly, still relies on fine-tuning a model on labeled
examples. Instruction-tuned large language models offer a different
route: Brown et al.~\cite{brown2020gpt3} showed that a sufficiently
capable model can perform a new task from a natural-language description
and a handful of demonstrations, with no weight update at all, which
makes prompting an attractive option whenever labeled data, compute, or
both are scarce.

That option comes with its own open question: which prompt does the
job. A prompt for sexism detection can assign the model a role, define
the linguistic markers that make language sexist, supply surrounding
context about the platform and content type, ask the model to reason
before answering, or show worked examples, each a real design choice
with no established default for this specific task. Related work on
prompt engineering generally, and on hate speech detection with LLMs
specifically (Sec.~\ref{sec:related}), suggests these choices can matter
a great deal, but the evidence for sexism detection in particular has so
far come from single models evaluated on a narrow slice of that design
space, without a non-LLM baseline or formal significance testing to say
whether an apparent improvement is real.

This paper closes that gap directly. We evaluate a curated grid of 18
prompt variants, combining five independently toggleable components
(role, aspects, context, chain-of-thought, and few-shot demonstrations)
in isolation and in combination, across six open-weight, instruction-tuned
models spanning five model families and a 3.8--9B parameter range, on the
full official EDOS test split. Every variant's effect against a
zero-component baseline is tested for statistical significance with two
independent tests, and every model's best training-free variant is
compared directly against a QLoRA-tuned version of the same model, a
fine-tuned encoder, and a classical TF-IDF and logistic regression
baseline trained on the same data. We further break results down by
sexism subtype and test three robustness axes prior work has not applied
to this task: the order in which prompt components appear, the number of
few-shot demonstrations shown, and which specific demonstrations are
drawn.

The results argue for a narrower and more specific claim than either
``prompt design matters'' or ``LLMs solve sexism detection'' taken
alone. Prompt design does produce large, statistically real effects
within training-free prompting: the large majority of this study's
variant comparisons reach significance under each of two independent
tests, and the single largest effect among this study's
significance-tested comparisons comes from prompt design alone. But
that entire space sits, for five of the six evaluated models, below
what a linear classifier trained on the same labeled data already
achieves, at a fraction of the computational cost; only one model
clears that bar on F1, at a steep compute premium, and a separate
confidence-based AUC comparison crowns a different model instead.
Along the way, reading raw model output rather than trusting aggregate
scores surfaced a chain-of-thought-triggered failure mode in one model
that splits cleanly into two distinct causes depending on prompt
composition, and showed that prompt-component order, unlike few-shot
demonstration order in prior work, has a comparatively narrow effect on
these models' behavior. Sec.~\ref{sec:limitations} and
Sec.~\ref{sec:conclusion} address how far these findings generalize
beyond the setting tested here.
